\begin{textsample}{POS Dim 3 – 2009\_c4 – Score 2.00 – 2009\_c4/enade\_2009\_est\_f2e70bb570.txt}  \label{ex:f3_pos_121}
Capítulo 4 Percepção sobre a Prova As percepções dos estudantes sobre a prova aplicada no ENADE/2009 na área de ESTATÍSTICA foram mensuradas por meio de nove itens que avaliaram desde o grau de dificuldade da prova até o tempo gasto para concluí -la.
As questões foram analisadas separando -se concluintes e ingressantes, e as percepções sobre a prova foram relacionadas com o desempenho dos estudantes e com a Região Geográfica de localização do curso6.
O desempenho dos estudantes dentro de cada Área/Subárea foi classificado em quartos ( 1/4 ) para os concluintes e para os ingressantes em separado ( ver Figura 1 para o esquema de representação gráfica destas grandezas ).
As notas finais foram colocadas em ordem ascendente e no primeiro quarto foram colocadas as 25 \% notas mais baixas, limitadas superiormente pelo percentil 25 que vem a ser o primeiro quartil.
No segundo quarto, as 25 \% seguintes, limitadas superiormente pelo percentil 50 que vem a ser a mediana.
No terceiro quarto, as 25 \% notas após as encontradas no segundo quarto, ou seja, as primeiras 25 \% acima da mediana, limitadas superiormente pelo percentil 75 que vem a ser o terceiro quartil.
No último quarto foram alocadas as 25 \% maiores notas.
Potencialmente, poderia existir uma relação entre a nota do estudante e a sua percepção sobre a prova.
Calculando as médias das respostas segundo os quartos de notas, seria possível reconhecer a existência ( ou não ) de algum gradiente nas respostas, por exemplo, uma melhor nota correspondendo a uma percepção menos negativa.
Este gradiente seria traduzido, por exemplo, no gráfico 4.1a, pela diminuição das áreas \textbf{laranja} e vermelha ( Difícil e Muito difícil ) com o aumento do desempenho ( indo do primeiro ao último quarto ). 6 Curso, neste relatório, refere -se à unidade de análise para o Conceito ENADE e é caracterizado pela combinação de IES, Área/Subárea ( quando pertinente ) e município de habilitação.
A medida tradicional de “ grau de discriminação ” comparou os resultados dos estudantes do primeiro quarto ( 25 \% piores estudantes ) aos do último quarto ( 25 \% melhores ).
Tal comparação, no gráfico disponibilizado, foi idêntica à comparação da primeira barra ( primeiro quarto ) com a última barra ( último quarto ) de cada grupo ingressante ou concluinte nos gráficos “ segundo quartos de nota final ” ( gráficos com números ímpares neste capítulo ).
Por outro lado, a disponibilização das informações segundo os quartos de nota permitiu verificar o comportamento de dois outros grupos intermediários que na análise somente dos extremos tem que ser suposto monotônico.
A seguir, serão apresentados os principais resultados relativos aos nove itens avaliados. 4.1 Grau de dificuldade da prova 4.1.1 Formação Geral Os Gráficos 4.1a, 4.1b e 4.1c apresentam a distribuição das respostas à questão “ Qual o grau de dificuldade desta prova na parte de Formação Geral? ”, segundo quartos de nota final e grupo de estudantes concluintes e ingressantes que participaram do ENADE/2009 na área de ESTATÍSTICA.
No total de concluintes, 24,0 \% declararam que a prova na parte da Formação Geral foi Difícil/Muito difícil, contra um valor de 21,1 para os ingressantes.
Tanto entre os concluintes, como entre os ingressantes, a proporção de indivíduos que classificaram esta parte como Difícil/Muito difícil diminui, grosso modo, com o desempenho do estudante.
Por outro lado, a proporção que declarou a parte Fácil/Muito fácil foi maior entre os concluintes ( 20,4 \% ) do que entre os ingressantes ( 16,4 \% ).
Não se observou uma variação monotônica significativa desta proporção com o aumento no desempenho, nem entre os concluintes, nem entre os ingressantes.
Entre os concluintes, 55,5 \% classificaram esta parte com um nível médio de dificuldade.
A proporção entre os ingressantes foi maior, 62,4 \%.
Entre os concluintes, esta proporção aumentou com o desempenho até o 3º quarto indo de 45,6 \% 1º quarto a 64,9 \% no 3º quarto caindo depois no último quarto para 53,0 \%.
Os Gráficos 4.2a e 4.2b apresentam a distribuição das respostas à questão “ Qual o grau de dificuldade desta prova na parte de Formação Geral? ”, segundo a Região onde o curso está situado para os estudantes concluintes e ingressantes que participaram do ENADE/2009 na área de ESTATÍSTICA.
Entre os concluintes, a maior proporção de estudantes que considerou a parte de Formação Geral da prova como de dificuldade média entre todas as Regiões brasileiras, foram os da Região Sul com 61,6 \% e entre os ingressantes, foram os estudantes da Região Nordeste com 67,2 \%.
A Região com as menores proporções desta resposta foram para os concluintes e ingressantes, a Centro-Oeste ( respectivamente 45,5 \% e 47,5 \% ). 4.1.2 Componente Específico Os Gráficos 4.3a, 4.3b e 4.3c apresentam a distribuição das respostas à questão “ Qual o grau de dificuldade desta prova na parte do Componente Específico? ”, segundo quartos de nota final e grupo de estudantes concluintes e ingressantes que participaram do ENADE/2009 na área de ESTATÍSTICA.
Em geral, para as desagregações consideradas, os estudantes classificaram o Componente Específico como mais difícil que a Formação Geral.
No grupo de concluintes, 36,5 \% declararam que a prova na parte do Componente Específico foi Difícil/Muito difícil, sendo um valor muito maior para os ingressantes, 57,6 \%.
Entre os concluintes, a proporção de indivíduos que classificaram esta parte como Muito difícil diminuiu com o desempenho do estudante.
A proporção que declara Difícil/Muito difícil entre os concluintes é, grosso modo, decrescente com o desempenho.
Entre os ingressantes essa proporção é razoavelmente estável em torno da média ( 57,6 \% ), tanto que a distância em pontos percentuais do maior valor para o menor valor é menor do que 5.
Entre os concluintes, aqueles \textbf{localizados} no segundo e no terceiro quarto não consideraram essa parte Muito fácil.
Entre os ingressantes, houve essa consideração nos três primeiros quartos de nota final.
A proporção que declarou a parte Fácil/Muito fácil foi maior entre os concluintes ( 6,7 \% ) do que entre os ingressantes ( 4,3 \% ).
Entre os concluintes a proporção dos que consideraram esta parte como Fácil/Muito Fácil não apresenta comportamento monotônico.
Um padrão decrescente é experimentado pelos ingressantes.
Entre os concluintes, 56,8 \% classificaram esta parte com um nível médio de dificuldade.
A proporção entre os ingressantes foi menor, 38,1 \%.
Entre os concluintes, esta proporção aumentou com o desempenho, grosso modo, indo de 44,8 \% no primeiro quarto de notas a 67,5 \% no último quarto.
Entre os ingressantes o comportamento é semelhante, indo de 35,4 \% a 42,2 \%.
Os Gráficos 4.4a e 4.4b apresentam a distribuição das respostas à questão “ Qual o grau de dificuldade desta prova na parte do Componente Específico? ”, segundo a Região onde o curso está situado para os estudantes concluintes e ingressantes que participaram do ENADE/2009 na área de ESTATÍSTICA.
Nenhum concluinte da Região Norte e da Sul declarou que esta parte foi muito fácil e entre os ingressantes a não declaração foi para a opção fácil na Região Norte.
Entre os concluintes, a maior proporção de estudantes que considerou a parte do Componente Específico da prova como de dificuldade média entre todas as Regiões brasileiras, foi na Região Centro-Oeste com 77,3 \%, já entre os ingressantes, foram os estudantes da Região Nordeste com 41,3 \%.
A Região com as menores proporções desta resposta foram respectivamente para os concluintes e ingressantes, a Norte ( 37,1 \% ) e a Centro-Oeste ( 28,3 \% ). 4.2 Extensão da prova em relação ao tempo total Os Gráficos 4.5a, 4.5b e 4.5c apresentam a distribuição das respostas à questão sobre a extensão da prova em relação ao tempo total, segundo quartos de nota final e grupo de estudantes concluintes e ingressantes que participaram do ENADE/2009 na área de ESTATÍSTICA.
Tanto para concluintes quanto para ingressantes, a maioria considerou a extensão da prova como adequada em relação ao tempo destinado à sua resolução, respectivamente 54,0 \% e 58,8 \%.
Entre os concluintes, esta proporção cresceu com o desempenho.
A proporção dos que consideraram a extensão da prova muito curta foi ligeiramente maior entre os concluintes ( 1,8 \% ) do que entre os ingressantes ( 1,5 \% ).
A proporção dos que consideraram a extensão da prova Curta/Muito curta ao tempo total foi menor entre os concluintes ( 4,7 \% ) do que entre os ingressantes ( 6,7 \% ).
Entre os concluintes, somente estudantes com desempenho no 3º grupo não declararam que a extensão da prova teria sido Muito curta com relação ao tempo total.
Entre os ingressantes, esta declaração aconteceu somente o último grupo.
A maior proporção de concluintes declarando que a prova foi Curta/Muito curta foi entre os estudantes do 2º quarto, 9,3 \%.
Entre os ingressantes, foi no 1º quarto, com 9,2 \%.
A proporção dos que consideraram a extensão da prova Muito longa foi maior entre os concluintes ( 17,7 \% ) do que entre ingressantes ( 10,4 \% ).
A proporção dos que consideraram a extensão da prova Longa foi menor entre os concluintes ( 23,6 \% ) do que entre os ingressantes ( 24,0 \% ).
Os Gráficos 4.6a e 4.6b apresentam a distribuição das respostas à questão sobre a extensão da prova em relação ao tempo total, segundo a Região onde o curso está situado para os estudantes concluintes e ingressantes que participaram do ENADE/2009 na área de ESTATÍSTICA.
Em todas as Regiões brasileiras, a maioria dos estudantes, tanto concluintes como ingressantes, perceberam a prova como adequada. 4.3 Compreensão dos Enunciados das Questões 4.3.1 Formação Geral Os Gráficos 4.7a, 4.7b e 4.7c apresentam a distribuição das respostas à questão “ Os enunciados das questões da prova na parte de Formação Geral, estavam claros e objetivos? ”, segundo quartos de nota final e grupo de estudantes concluintes e ingressantes que participaram do ENADE/2009 na área de ESTATÍSTICA.
Com relação aos enunciados das questões da parte de Formação Geral, as opiniões foram positivas : a grande maioria dos concluintes ( 73,5 \% ) e ingressantes ( 78,9 \% ) consideram claros e objetivos os enunciados de todas ou da maioria das questões.
Nem entre os concluintes e nem entre ingressantes distingue -se uma tendência no comportamento em função do desempenho.
Os níveis de resposta informando que Não, nenhum dos enunciados das questões da parte de Formação Geral era claro e objetivo foram muito baixos, tanto para concluintes ( 1,3 \% ) como para ingressantes ( 2,1 \% ).
A proporção de estudantes com esse tipo de declaração decresce com o desempenho tanto para concluintes como para ingressantes : ingressantes ( de 3,8 \% no primeiro quarto para 0,9 \% no último quarto ) e concluintes ( de 3,0 \% no primeiro quarto para 0,0 \% no último quarto ).
Os Gráficos 4.8a e 4.8b apresentam a distribuição das respostas à questão “ Os enunciados das questões da prova na parte de Formação Geral, estavam claros e objetivos? ”, segundo a Região onde o curso está situado para os estudantes concluintes e ingressantes que participaram do ENADE/2009 na área de ESTATÍSTICA.
Em termos regionais, as proporções de ingressantes e de concluintes que consideraram claros e objetivos os enunciados de todas ou da maioria das questões variaram no máximo 12,4 pontos percentuais dentro de um grupo.
Entre os concluintes, os valores foram de 69,4 \% no Sul a 81,8 \% no Centro-Oeste.
Já entre os ingressantes estas proporções variaram de 76,7 \% no Nordeste a 83,9 \% no Sul.
Com relação às respostas Poucos e Não, nenhum tanto para concluintes como para ingressantes, as maiores proporções ocorreram nas Regiões Norte e Nordeste, respectivamente 14,3 \% e 14,1 \%. 4.3.2 Componente Específico Os Gráficos 4.9a, 4.9b e 4.9c apresentam a distribuição das respostas à questão “ Os enunciados das questões da prova na parte de formação específica estavam claros e objetivos? ”, segundo quartos de nota final e grupo de estudantes concluintes e ingressantes que participaram do ENADE/2009 na área de ESTATÍSTICA.
Com relação aos enunciados das questões da parte do Componente Específico, as opiniões foram também positivas : a grande maioria dos concluintes ( 73,5 \% ) e ingressantes ( 71,2 \% ) consideraram claros e objetivos os enunciados de todas ou da maioria das questões.
A combinação de todas ou da maioria das questões é, grosso modo, crescente com o desempenho, indo de 71,7 \% para os com pior desempenho a 77,5 \% para o último quarto entre os concluintes.
Entre os ingressantes, esta proporção cresce com o desempenho indo de 64,0 \% para os com pior desempenho a 80,0 \% para o último quarto.
O número de estudantes que declararam que não, nenhum dos enunciados das questões da parte do Componente Específico era claro e objetivo, foi maior do que na parte de Formação Geral : 0,8 \% entre os concluintes e 1,7 \% entre os ingressantes.
Os que responderam Poucos são 8,1 \% entre os concluintes e 8,7 \% entre os ingressantes.
A proporção de estudantes que respondeu que cerca de metade dos enunciados da parte do Componente Específico era clara e objetiva decresceu com o desempenho entre os ingressantes, caindo de 20,8 \% a 14,5 \%.
Os Gráficos 4.10a e 4.10b apresentam a distribuição das respostas à questão “ Os enunciados das questões da prova na parte de formação específica estavam claros e objetivos? ”, segundo a Região onde o curso está situado para os estudantes concluintes e ingressantes que participaram do ENADE/2009 na área de ESTATÍSTICA.
Em termos regionais, as proporções de ingressantes e de concluintes que consideraram claros e objetivos os enunciados de todas ou da maioria das questões variaram mais entre os ingressantes ( no máximo 21,4 pontos percentuais ) do que entre os concluintes ( no máximo 16,6 pontos percentuais ).
Entre os concluintes, os valores foram de 63,6 \% no Centro-Oeste a 80,2 \% no Sudeste.
Entre os ingressantes, os valores foram de 67,8 \% na Região Nordeste a 78,0 \% na Região Sul.
Com relação às respostas “ poucos ” e “ não, nenhum ” para concluintes as maiores proporções ocorreram, respectivamente, na Região Norte com 17,7 \% e na Região Nordeste com 15,6 \%. 4.4 Suficiência das informações/instruções fornecidas Os Gráficos 4.11a, 4.11b e 4.11c apresentam a distribuição das respostas à questão “ As informações/instruções fornecidas para a resolução das questões foram suficientes para resolvê -las? ”, segundo quartos de nota final e grupo de estudantes concluintes e ingressantes que participaram do ENADE/2009 na área de ESTATÍSTICA.
Quando perguntados sobre as informações/instruções fornecidas para a resolução das questões, quase 78,0 \% dos ingressantes e 82,7 \% dos concluintes de todo o Brasil responderam que estas eram ou excessivas ou suficientes em todas ou na maioria das questões, demonstrando uma avaliação positiva da prova.
A proporção de estudantes com a opinião supracitada foi crescente com o desempenho para concluintes, com valores indo de 66,2 \% para os de pior desempenho a 92,0 \% para os de melhor desempenho.
Para os ingressantes esta proporção foi crescente somente até o terceiro quarto de desempenho, com valores indo de 74,1 \% para os de pior desempenho a 84,1 \% no terceiro quarto, seguido de uma queda para 79,1 \% no último quarto.
O complemento desta opinião, ou seja, os estudantes que responderam Sim, somente em algumas ou Não, em nenhuma delas, apresentou uma proporção decrescente com o desempenho para os concluintes.
Os Gráficos 4.12a e 4.12b apresentam a distribuição das respostas à questão “ As informações/instruções fornecidas para a resolução das questões foram suficientes para resolvê -las? ”, segundo a Região onde o curso está situado para os estudantes concluintes e ingressantes que participaram do ENADE/2009 na área de ESTATÍSTICA. 4.5 Dificuldade encontrada ao responder à prova Os Gráficos 4.13a, 4.13b e 4.13c apresentam a distribuição das respostas à questão “ Você se deparou com alguma dificuldade ao responder à prova? ”, segundo quartos de nota final e grupo de estudantes concluintes e ingressantes que participaram do ENADE/2009 na área de ESTATÍSTICA.
Perguntados sobre as dificuldades com as quais se depararam ao responder à prova, 62,4 \% dos ingressantes apontaram como principal dificuldade o desconhecimento do conteúdo, já os concluintes apontaram como principal dificuldade a forma diferente de abordagem, 36,3 \%.
A forma diferente de abordagem do conteúdo foi considerada como maior dificuldade por 15,6 \% dos ingressantes, enquanto que a falta de motivação foi a principal dificuldade para 32,9 \% dos concluintes e para 14,4 \% dos ingressantes.
Um percentual maior de concluintes ( 10,4 \% ) do que de ingressantes ( 3,7 \% ) responderam não ter tido nenhuma dificuldade para responder à prova.
Os Gráficos 4.14a e 4.14b apresentam a distribuição das respostas à questão “ Você se deparou com alguma dificuldade ao responder à prova? ”, segundo a Região onde o curso está situado para os estudantes concluintes e ingressantes que participaram do ENADE/2009 na área de ESTATÍSTICA.
Em relação às Regiões brasileiras, o percentual de ingressantes que apontaram o desconhecimento do conteúdo como principal dificuldade ao responder à prova variou de 55,7 \% na Região Centro-Oeste, a 75,0 \% na Norte.
Entre os concluintes, tal percentual variou de 5,5 \% na Região Sul a 14,3 \%, nas Regiões Sudeste e Norte.
A forma diferente de abordagem de conteúdo foi citada como maior dificuldade por 18,2 \% dos concluintes na Região Centro-Oeste chegando a 45,7 \% nas Regiões Norte e Nordeste.
Já entre os ingressantes, os valores foram de 11,5 \% na Região Centro-Oeste a 19,2 \% na Região Nordeste.
O percentual de concluintes que citaram a falta de motivação como principal dificuldade variou entre os concluintes de 14,3 \% na Região Norte a 50,0 \% na Região Centro-Oeste.
Entre os ingressantes, as estatísticas correspondentes foram de 3,8 \% na Região Norte a 21,3 \%, na Região Centro-Oeste.
Os que declararam que não tiveram qualquer tipo de dificuldade variaram, entre os concluintes, de 6,5 \% na Região Nordeste a 20,0 \% na Região Norte.
Entre os ingressantes, este percentual variou de 0,0 \% na Região Norte a 6,6 \% na Região Centro-Oeste. 4.6 Percepção sobre os conteúdos das questões objetivas da prova Os Gráficos 4.15a, 4.15b e 4.15c apresentam a distribuição das respostas à questão “ Considerando apenas as questões objetivas da prova, você percebeu que :... ”, segundo quartos de nota final e grupo de estudantes concluintes e ingressantes que participaram do ENADE/2009 na área de ESTATÍSTICA.
As respostas possíveis para esta questão estavam relacionadas à cobertura do conteúdo das questões objetivas durante o curso.
A proporção de ingressantes ( 65,7 \% ) que declarou que não estudou ainda a maioria desses conteúdos foi bem maior do que entre os concluintes ( 3,1 \% ).
Esta proporção, entre os concluintes, foi decrescente como função do desempenho, embora haja um pequeno aumento no segundo quarto de notas.
Entre os concluintes, essa proporção foi decrescente com o desempenho a partir do segundo quarto.
Por outro lado, a proporção de estudantes que declarou que estudou e aprendeu todos ou a maioria desses conteúdos, entre os ingressantes foi de 14,9 \% e entre os concluintes 65,4 \%.
Esta proporção, entre concluintes, foi crescente como função do desempenho.
A proporção de estudantes que declarou ter estudado a maioria desses conteúdos, mas não os ter aprendido foi bem maior entre os concluintes ( 5,5 \% ) do que os ingressantes ( 22,7 \% ).
Os Gráficos 4.16a e 4.16b apresentam a distribuição das respostas à questão “ Considerando apenas as questões objetivas da prova, você percebeu que :... ”, segundo a Região onde o curso está situado para os estudantes concluintes e ingressantes que participaram do ENADE/2009 na área de ESTATÍSTICA.
Em relação às Regiões brasileiras, concluintes que afirmaram não ter estudado ainda a maioria dos conteúdos só não ocorreram na Região Norte e variaram de 0,0 \% nesta Região a 9,8 \% na Nordeste.
Entre os ingressantes, os percentuais variaram de 48,3 \% na Região Centro-Oeste a 77,4 \%, na Norte.
Em todas as Regiões, a maioria dos concluintes afirmou ter estudado e aprendido todos ou muitos dos conteúdos, com percentuais que variaram de 50,0 \% na Região Nordeste a 81,8 \%, na Centro-Oeste.
Para os ingressantes, tal percentual variou de 5,7 \%, na Região Norte, a 21,6 \%, na Centro-Oeste. 4.7 Tempo gasto para concluir a prova Os Gráficos 4.17a, 4.17b e 4.17c apresentam a distribuição das respostas à questão “ Qual foi o tempo gasto para concluir a prova? ”, segundo quartos de nota final e grupo de concluintes e ingressantes que participaram do ENADE/2009 na área de ESTATÍSTICA.
Quando perguntados quanto ao tempo de conclusão da prova, a maioria dos concluintes indicou entre duas e três horas ( 37,9 \% ) seguido de entre três e quatro horas ( 28,6 \% ) e entre uma e duas horas ( 22,4 \% ).
Entre os ingressantes, o comportamento desta questão seguiu outra ordem, pois a maioria indicou entre uma e duas horas ( 39,5 \% ), seguido por entre duas e três horas ( 34,4 \% ) e entre três e quatro horas ( 15,3 \% ).
A proporção daqueles que declararam ter usado as quatro horas, mas não terem conseguido terminar foi baixa : 1,7 \% entre os ingressantes e 6,7 \% entre os concluintes.
Nem entre os concluintes, nem entre os ingressantes percebeu -se uma tendência monotônica do tempo necessário para concluir a prova como função do desempenho.
Entre os concluintes, a distribuição do tempo mostra que os estudantes do terceiro quarto levaram menos tempo e os do primeiro quarto mais tempo na confecção da prova.
Entre os ingressantes, os mesmo se repete alterando apenas os quartos, pois os do terceiro ocuparam mais tempo e os do segundo menos tempo na confecção da prova.
Os Gráficos 4.18a e 4.18b apresentam a distribuição das respostas à questão “ Qual foi o tempo gasto para concluir a prova? ”, segundo a Região onde o curso está situado para os estudantes concluintes e ingressantes que participaram do ENADE/2009 na área de ESTATÍSTICA.
A proporção de concluintes que gastaram entre uma e duas horas foi maior na Região Sul ( 34,2 \% ) e menor na Região Norte ( 8,6 \% ).
Para ingressantes, o maior percentual para este tempo ocorreu na Região Sul ( 46,1 \% ) e o menor na Região Norte ( 30,2 \% ).
Considerando -se a distribuição cumulativa, entre os concluintes, os estudantes da Região Norte foram os que levaram mais tempo para concluir a prova, por oposição aos estudantes da Região Sul que levaram menos tempo.
Entre os ingressantes, as Regiões correspondentes foram a Nordeste, com tempo mais longo e a Sudeste, com tempo mais curto.

% matched lemmas: laranja, localizar
\end{textsample}
